% Generated from locale/tr/content/sets-functions-relations/size-of-sets/size-of-sets-complete.tex; do not hand-edit.


\allmoduleschapter{sfr}{siz}{Kümelerin Büyüklüğü}

\begin{editorial}
Bu bölüm numaralandırmaları, sayılabilirliği ve sayılamazlığı ele alır.
Bazı kesimlerin iki sürümü vardır: daha temel sürüm, numaralandırmaları
listeler veya $\PosInt$'tan örten fonksiyonlar olarak ele alır; daha soyut
sürüm ise numaralandırmaları $\Nat$ ile bire bir örten eşlemeler olarak
tanımlar.
\end{editorial}

\allmodulesimport[section]{OLP-0028}

\allmodulesimport[section]{OLP-0029}

\allmodulesimport[section]{OLP-0030}

\allmodulesimport[section]{OLP-0031}

\allmodulesimport[section]{OLP-0032}

\allmodulesimport[section]{OLP-0033}

\allmodulesimport[section]{OLP-0034}

\allmodulesimport[section]{OLP-0035}

\allmodulesimport[section]{OLP-0036}

\allmodulesimport[section]{OLP-0037}

\begin{editorial}
Aşağıdaki \olref[sfr][siz][enm-alt]{sec},
\olref[sfr][siz][nen-alt]{sec} ve \olref[sfr][siz][red-alt]{sec}, Tim
Button'ın \emph{Open Set Theory} metninde kullanılmak üzere hazırladığı,
sırasıyla \olref[sfr][siz][enm]{sec}, \olref[sfr][siz][nen]{sec} ve
\olref[sfr][siz][red]{sec} kesimlerinin alternatif sürümleridir. Bunlar biraz
daha ileri düzeydedir ve küme teorisi bağlamına daha uygun farklı bir
sayılabilirlik tanımı kullanır (yani listelenebilirlik veya $\PosInt$'tan örten
bir fonksiyonun görüntü kümesi olmak yerine $\Nat$ ya da onun bir başlangıç
parçasıyla bire bir örten eşlenebilirlik).
\end{editorial}

\allmodulesimport[section]{OLP-0038}
\allmodulesimport[section]{OLP-0039}
\allmodulesimport[section]{OLP-0040}

\OLEndChapterHook

